\begin{table}[H]
\centering
\caption{RGARCH-CARR-SK 拟合准则}
\label{tab:rgarch-fit}
\small
\begin{threeparttable}
\begin{tabular}{lrrr}
\toprule
realized measure & LLK & AIC & BIC \\
\midrule
RV & -169.822 & 365.64 & 418.62 \\
RBV & -198.783 & 423.57 & 476.54 \\
MedRV & -203.959 & 433.92 & 486.90 \\
RMAD & -553.492 & 1,133 & 1,186 \\
\bottomrule
\end{tabular}
\begin{tablenotes}[flushleft]
\footnotesize
\item 注：LLK、AIC 和 BIC 用于比较不同 realized pressure measure 下的样本内拟合表现。
\end{tablenotes}
\end{threeparttable}
\end{table}

\begin{table}[H]
\centering
\caption{RGARCH-CARR-SK 样本外损失}
\label{tab:rgarch-loss}
\small
\begin{threeparttable}
\begin{tabular}{lrrr}
\toprule
realized measure & R2LOG(validation) & R2LOG(test) & MAE(test) \\
\midrule
RV & 15.605 & 15.396 & 53.392 \\
RBV & 15.578 & 15.266 & 53.420 \\
MedRV & 16.855 & 16.693 & 53.378 \\
RMAD & 0.020 & 0.030 & 0.131 \\
\bottomrule
\end{tabular}
\begin{tablenotes}[flushleft]
\footnotesize
\item 注：R2LOG 为损失指标，数值越低越好。RMAD 的 R2LOG 明显低于其他测度，但该差异与测度尺度有关，因此正文使用表格报告而不使用线性尺度柱状图。
\end{tablenotes}
\end{threeparttable}
\end{table}
